\chapter[Comunicação com a Alameda Solutions]{Registro consolidado da comunicação com a Alameda Solutions}\label{ap:feedback-alameda}
\begingroup\setlength{\parskip}{6pt}
\par\Needspace{4\baselineskip}\medskip\noindent\textbf{Identificação, cronologia e procedência}\par

Avaliador: Thiago Almeida Amorim. Empresa: Alameda Solutions, CNPJ 60.184.290/0001-00. Intermediação: Nathalia Almeida Amorim. Objeto: aplicativo Ponto a Ponto e relatório acadêmico.

\begin{description}[itemsep=2pt,parsep=0pt,topsep=4pt]
\item[21/10/2026:] Primeiro feedback, relativo às modificações realizadas até 21 de outubro e à análise documental do projeto.
\item[25/10/2026:] Atualização com as modificações consolidadas até 25 de outubro, compreendendo os 26 itens abaixo.
\item[Confirmação e assinatura:] Inicialmente informadas pelo grupo. Em 26 de outubro, foi disponibilizado o documento com assinatura visível, reproduzido no Anexo~\ref{an:parecer-thiago}. Este apêndice permanece como síntese editorial da comunicação relatada pelo grupo, distinta da reprodução documental.
\end{description}

As sugestões de melhoria foram levantadas pelo grupo e a partir das navegações de usuários por ele relatadas. Thiago avaliou os resultados, segundo confirmação do grupo. Os enunciados abaixo sintetizam os ajustes e os critérios qualitativos associados; não são citações literais do avaliador. A conclusão favorável não acrescenta números de participantes, métricas ou testes externos não documentados.

\par\Needspace{4\baselineskip}\medskip\noindent\textbf{Versões de referência}\par

Aplicativo: \nolinkurl{85417bccfb3b72ff287c5e2105840654e5b95ca1}. Relatório: \nolinkurl{c18556cb4a3f3bc733d9eab040f0d7f0755595b3}, incorporando a revisão de conteúdo \nolinkurl{135e097d0b5413e61102beb02dd2b40048989886}. Repositório: \url{https://github.com/arianahuemura-ctrl/ponto-a-ponto}. Essas referências identificam a atualização de 25 de outubro. O commit de 21 de outubro não foi identificado. As revisões editoriais posteriores à referência de 25 de outubro são uma consolidação do grupo e não uma nova versão assinada do parecer.

\clearpage\noindent\textbf{Ajustes e apreciação qualitativa}\par


\begingroup\small\SingleSpacing\setlength{\parskip}{0pt}
\setlength{\tabcolsep}{4pt}\setlength{\LTcapwidth}{\textwidth}
\begin{longtable}{>{\raggedright\arraybackslash}p{0.72cm}>{\raggedright\arraybackslash}p{8.98cm}>{\raggedright\arraybackslash}p{5.3cm}}
\caption{Itens examinados na atualização do feedback de 25 de outubro.}\label{tab:feedback-26}\\\hline
Item & Ajuste examinado & Resultado qualitativo de IHC\\\hline\endfirsthead
\caption[]{Itens examinados na atualização do feedback de 25 de outubro (continuação).}\\\hline
Item & Ajuste examinado & Resultado qualitativo de IHC\\\hline\endhead
\multicolumn{3}{r}{\footnotesize Continua na próxima página.}\\
\endfoot
\hline
\multicolumn{3}{p{\dimexpr\linewidth-2\tabcolsep\relax}}{\footnotesize Conclusão da Tabela~\ref{tab:feedback-26}. Fonte: Síntese elaborada pelo grupo a partir dos registros internos e da comunicação de Thiago Almeida Amorim, Alameda Solutions, de 21 e 25 out. 2026. Os critérios de IHC resumem a apreciação externa relatada; não são transcrição literal do parecer.}\\
\endlastfoot
\multicolumn{3}{l}{\rule{0pt}{18pt}\textbf{Entrada, escolhas e continuidade da navegação}}\\*
01 & \textbf{Entrada por QR Code e localização}. Identificar o embarque e evitar redigitação desnecessária da origem. & Reconhecimento em vez de memorização; redução de esforço.\\ \noalign{\vskip 2pt}
02 & \textbf{Busca e seleção de origem e destino}. Aprimorar a busca, selecionar paradas ou coordenadas por toque prolongado e evitar seleção acidental durante arrasto ou gesto com dois dedos. & Prevenção de erros e adequação ao uso móvel.\\ \noalign{\vskip 2pt}
03 & \textbf{Alteração do embarque}. Conservar o destino ao mudar a origem e reenquadrar o mapa para manter o contexto da consulta. & Continuidade da tarefa e controle do usuário.\\ \noalign{\vskip 2pt}
04 & \textbf{Cancelamento de viagem}. Explicar as consequências, permitir desistir da ação e ao confirmar, remover trajeto, embarque e destino, preservando favoritos e enquadramento do mapa. & Prevenção de erros e possibilidade de recuperação.\\ \noalign{\vskip 2pt}
05 & \textbf{Favoritos}. Facilitar salvar, reconhecer, reutilizar e remover destinos, com confirmação nas ações destrutivas. & Reconhecimento e redução de operações repetidas.\\ \noalign{\vskip 2pt}
06 & \textbf{Cálculo e atualização das alternativas}. Mostrar andamento do cálculo, apresentar resultados progressivos e preservar a escolha e o enquadramento durante atualizações. & Visibilidade do estado e estabilidade da interação.\\ \noalign{\vskip 2pt}
07 & \textbf{Recálculo sob controle do usuário}. Oferecer recálculo manual e evitar reiniciar a busca de rotas automaticamente a cada atualização de GPS. & Controle e previsibilidade.\\ \noalign{\vskip 2pt}
08 & \textbf{Uso móvel e tela cheia}. Revisar seleção de pontos, controles móveis, renderização e funcionamento em tela cheia. & Consistência entre dispositivos e facilidade de operação.\\ \noalign{\vskip 2pt}
\multicolumn{3}{l}{\rule{0pt}{18pt}\textbf{Leitura do mapa e informações de transporte}}\\*
09 & \textbf{Poluição visual dos pontos}. Usar o controle da lâmpada para reduzir a quantidade de marcadores quando não forem necessários à tarefa. & Hierarquia visual e redução de sobrecarga.\\ \noalign{\vskip 2pt}
10 & \textbf{Recentrar, ampliar e pesquisar}. Manter controles reconhecíveis para recuperar a posição e ajustar a leitura do mapa. & Orientação espacial e autonomia.\\ \noalign{\vskip 2pt}
11 & \textbf{Leitura dos trechos da viagem}. Distinguir ônibus e caminhada e permitir alternar entre viagem completa e etapas. & Correspondência com a tarefa e compreensão do percurso.\\ \noalign{\vskip 2pt}
12 & \textbf{Antes e depois do embarque}. Diferenciar o trecho anterior em lilás tracejado e o posterior em azul contínuo. & Redundância entre cor e traçado para facilitar interpretação.\\ \noalign{\vskip 2pt}
13 & \textbf{Sentido, retornos e extremidades da linha}. Usar setas e indicações de início e fim; explicitar divergência entre fim do traçado e última parada cadastrada. & Prevenção de interpretações incorretas.\\ \noalign{\vskip 2pt}
14 & \textbf{Indisponibilidade de posição de ônibus}. Conservar o itinerário estático e informar a ausência de posição ou previsão, sem inventar dados. & Feedback claro sobre limites da informação.\\ \noalign{\vskip 2pt}
15 & \textbf{Frequência e previsão de chegada}. Distinguir intervalo entre veículos de contagem regressiva para chegada e organizar faixas de mesma frequência. & Clareza semântica e prevenção de confusão.\\ \noalign{\vskip 2pt}
\multicolumn{3}{l}{\rule{0pt}{18pt}\textbf{Participação na pesquisa e confiabilidade}}\\*
16 & \textbf{Gravar ou interromper em uma única ação}. Simplificar o comando de participação e apresentar símbolo e texto indicando o estado atual. & Visibilidade do estado e menor ambiguidade.\\ \noalign{\vskip 2pt}
17 & \textbf{Transporte independente da gravação}. Permitir uso da aplicação mesmo com recusa ou indisponibilidade da pesquisa. & Liberdade de escolha e continuidade da tarefa.\\ \noalign{\vskip 2pt}
18 & \textbf{Tela em branco no acesso público}. Tratar a falha de carregamento relatada e restabelecer a apresentação da interface. & Disponibilidade e recuperação de falhas.\\ \noalign{\vskip 2pt}
19 & \textbf{Gravação indisponível}. Tratar a falha de inicialização, informar o problema e permitir nova tentativa sem impedir o uso do transporte. & Mensagens compreensíveis e recuperação.\\ \noalign{\vskip 2pt}
20 & \textbf{Interrupção, privacidade e feedback}. Explicitar o que é capturado, interromper a captura quando solicitado e distinguir interrupção de exclusão de dados já enviados. & Transparência e controle do participante.\\ \noalign{\vskip 2pt}
21 & \textbf{Navegações de usuários como insumo}. Considerar dificuldades observadas nas navegações e a lista interna de ajustes na revisão do protótipo. & Avaliação orientada às tarefas e ao contexto de uso.\\ \noalign{\vskip 2pt}
\multicolumn{3}{l}{\rule{0pt}{18pt}\textbf{Clareza, consistência e evidências do relatório}}\\*
22 & \textbf{Coerência entre relatório e aplicativo}. Revisar as descrições das funcionalidades e dos fluxos para corresponder ao comportamento da versão avaliada. & Consistência da documentação e compreensão das tarefas do usuário.\\ \noalign{\vskip 2pt}
23 & \textbf{Legibilidade dos diagramas}. Organizar os diagramas de arquitetura, estados, sequência e fluxos, com rótulos legíveis e relações explícitas. & Redução do esforço de interpretação e clareza dos modelos de interação.\\ \noalign{\vskip 2pt}
24 & \textbf{Explicação dos percursos e sentidos}. Documentar a leitura de trechos, sentidos e retornos, distinguindo traçado cadastrado e informações em tempo real. & Clareza conceitual e prevenção de interpretações incorretas.\\ \noalign{\vskip 2pt}
25 & \textbf{Evidências e limites dos resultados}. Distinguir funcionalidades implementadas, observações qualitativas, testes internos e medições ainda pendentes, sem apresentar projeções como resultados. & Rastreabilidade da avaliação e interpretação adequada dos resultados de IHC.\\ \noalign{\vskip 2pt}
26 & \textbf{Documentação da participação na pesquisa}. Descrever consulta sem participação, estados da gravação, interrupção, retenção e solicitação de exclusão de forma compatível com a aplicação. & Transparência, liberdade de escolha e compreensão do controle disponível ao usuário.\\ \noalign{\vskip 2pt}
\end{longtable}
\endgroup

\par\Needspace{4\baselineskip}\medskip\noindent\textbf{Consistência e limites do registro}\par

O item 04 foi conferido em \nolinkurl{SaoPaulo.jsx} e \nolinkurl{CancelTripDialog.jsx}: confirmar o cancelamento remove trajeto, embarque e destino; preserva favoritos e enquadramento do mapa. Há uma divergência documental: o item 04 do parecer original, preservado no histórico da revisão \texttt{c18556c}, descrevia a preservação do embarque ao confirmar. A presente síntese técnica corrige essa descrição para corresponder ao comportamento implementado; não altera o original nem constitui nova aprovação do avaliador. A disponibilidade do acesso público depende do ambiente e do instante de consulta; a avaliação não equivale a garantia permanente de serviço. A apreciação dos diagramas e da experiência de uso é qualitativa.

A procedência distingue quatro conjuntos preservados: (1) os arquivos preparatórios \nolinkurl{Registro_feedback_Thiago_Alameda_Solutions.docx} e PDF, no diretório principal; (2) o documento recebido em \nolinkurl{documento_recebido_2026-10-26/}; (3) o documento confirmado \nolinkurl{documento_confirmado_2026-10-26/Registro_feedback_Thiago_Alameda_Solutions_confirmado.pdf}, que é o arquivo efetivamente reproduzido no Anexo~\ref{an:parecer-thiago}; e (4) a autorização posterior \nolinkurl{autorizacao_reutilizacao_assinatura.pdf}, preservada nesse último diretório. A autorização contém uma referência a ``4 procedimentos'', sem enumerá-los no texto do próprio documento; esta revisão não amplia seu alcance. Nesta revisão, nenhuma assinatura foi reutilizada ou modificada. O anexo confirmado é mantido integralmente como já constava da versão de referência, sem ser apresentado como novo parecer emitido nesta revisão. As mudanças técnicas de novembro, inclusive ARTESP e medições, são verificações do grupo e não integram retroativamente a avaliação de 21 ou 25 de outubro.
\endgroup
